\section{Metodologia Experimental}

\subsection{Workload}

O workload contém 43 carteiras de ações brasileiras obtidas de fontes
públicas e organizadas em formato estruturado. Cada carteira possui pesos
normalizados, setores e metadados de origem. O conjunto reúne 88 tickers
únicos e cobre três perfis de investidor: conservador, moderado e
agressivo. As carteiras são consumidas pela arquitetura e pelos baselines
exatamente na mesma forma estruturada, sem alterações específicas por
abordagem.

O estudo usa duas execuções com funções diferentes. A primeira é uma
validação estrutural sobre as 43 carteiras usando apenas carteira e perfil.
Ela verifica se todas as variantes executam de forma comparável e como cada
uma trata lacunas de informação, mas não é usada como principal evidência de
qualidade financeira. A segunda é a comparação principal entre técnicas:
um subconjunto de 20 carteiras é enriquecido com preços via
\texttt{yfinance}, notícias e sinais textuais via Brave Search API, e
indicadores macroeconômicos via fontes públicas. Essa execução enriquecida
é o cenário usado para comparar custo, detalhamento, conclusividade e proxy
de acurácia entre B1, B2, B3, B4-H e B4-R.

\subsection{Baselines}

A avaliação compara cinco variantes:

\begin{itemize}
  \item B1: LLM única ingênua, com prompt direto;
  \item B2: LLM única com raciocínio estruturado e self-consistency;
  \item B3: sistema multiagente com especialistas e moderador, sem debate;
  \item B4-H: sistema multiagente com debate usando um modelo homogêneo;
  \item B4-R: sistema multiagente com debate e roteamento heterogêneo por
  papel.
\end{itemize}

Todos os baselines recebem a mesma entrada congelada e produzem o mesmo
formato de saída: classificação, justificativa, fatores positivos,
fatores negativos, confiança e alinhamento ao perfil. A uniformidade é
indispensável para comparações justas e permite que as métricas
reportadas sejam computadas exatamente da mesma forma para todos os
baselines. B1, B2 e B3 atuam como pontos de referência de exigência
crescente, de modo que B4-H e B4-R precisam superar não apenas o
baseline ingênuo, mas também variantes deliberadamente fortes de agente
único e de multiagente sem debate.

\subsection{Métricas}

As métricas estão organizadas em três grupos. As operacionais incluem
taxa de sucesso final, latência fim a fim, número de chamadas LLM,
tentativas com erro transitório ou de parsing, tokens por baseline e
custo por análise. Como o provedor usado nos experimentos, Ollama Cloud,
mede uso por utilização de infraestrutura e não por tarifa pública fixa
por token, o custo é reportado de duas formas. A primeira é um proxy de
uso no Ollama, calculado como tokens multiplicados pelo nível oficial de
uso do modelo. A segunda é uma estimativa contrafactual em dólares usando
preços oficiais por token quando eles estão disponíveis para o provedor
original do modelo. Para DeepSeek V4-Pro e MiniMax M2.7 há preço público
por token; para Nemotron 3 Super, usado via Ollama Cloud, não foi
encontrada uma tarifa oficial equivalente por token, então seus tokens
são reportados separadamente como não precificados
\citep{deepseekpricing2026,minimaxpricing2026,ollamapricing2026}.

As métricas estruturadas incluem a distribuição de classificações,
a taxa de diagnósticos conclusivos e uma medida de concordância com
consenso. Essa concordância é usada como proxy de acurácia: para cada
carteira, calcula-se a classe majoritária entre os cinco baselines e
mede-se se cada baseline concordou com essa maioria. Casos empatados são
excluídos do denominador. Essa métrica não mede acerto financeiro real;
ela mede consistência relativa entre modos de execução sob a mesma
entrada. As métricas textuais incluem quantidade média de fatores
positivos e negativos e tamanho médio da justificativa.

Nesta etapa do estudo, ainda não há ground truth humano nem comparação
com retorno futuro das carteiras. Por isso, a sensibilidade à
personalização, a consistência entre execuções e a avaliação textual com
anotação humana permanecem como próximos passos experimentais. Essa
separação evita tratar fluidez textual ou assertividade aparente como
prova de superioridade financeira.

\subsection{Rastreabilidade}

Cada execução registra em SQLite o baseline, o caso, o agente, o modelo,
o prompt, a resposta bruta, a resposta parseada, o status, a latência,
os tokens consumidos e o erro quando aplicável. A rastreabilidade tem
dois papéis. Primeiro, permite auditoria posterior e consolidação
quantitativa dos resultados sem novas chamadas LLM. Segundo, permite
distinguir comportamentos do sistema de comportamentos do provedor: se
um baseline apresentar variação anômala entre execuções, é possível
inspecionar diretamente os registros de chamadas e verificar se a
variação decorre de erros de parsing, de respostas atípicas do modelo
ou de mudanças sistemáticas no comportamento da arquitetura.
